% ============================================================
\chapter{Background and Related Work}
\label{chap:background}
% ============================================================

Rather than survey the two literatures this project sits between (portfolio
construction and causal discovery), this chapter analyses in depth the two
works it directly builds on, one per theme. L\'opez de Prado's Hierarchical
Risk Parity and backtest-rigour agenda (Section~\ref{sec:hrp-theme}) supplies
both the allocator whose symmetric clustering step C2 targets and the
inferential standard C3 adopts. Howard, Lohre and Mudde's causal networks on
the S\&P~500 (Section~\ref{sec:howard-theme}) are the closest causal-finance
work, and stop exactly one step short of the question this thesis asks.
Section~\ref{sec:cd-tools}
compresses the causal-discovery toolkit into the facts the pipeline relies
on, and Section~\ref{sec:gap} states the gap.

\section{Theme I: hierarchical allocation and honest backtests
(L\'opez de Prado)}
\label{sec:hrp-theme}

Mean--variance optimisation \cite{markowitz1952portfolio} requires inverting a
covariance matrix that, for a 100-asset universe estimated on a one-year
window, is noisy and near-singular; the optimiser amplifies the estimation
error and produces concentrated, unstable weights. Hierarchical Risk Parity
\cite{lopezdeprado2016hrp} avoids the inversion with a three-step scheme.
(i) Clustering: convert the correlation matrix $\rho$ into the
distance $d_{ij} = \sqrt{\tfrac{1}{2}(1-\rho_{ij})}$ and build a
single-linkage dendrogram. (ii) Quasi-diagonalisation: reorder assets
by the dendrogram's leaf order, concentrating large covariances near the
diagonal. (iii) Recursive bisection: split the ordered list in half,
allocate between halves in inverse proportion to their cluster variances, and
recurse. Because no matrix is inverted, HRP is robust to ill-conditioning and
beats mean--variance and inverse-variance portfolios out of sample in
the paper's Monte Carlo experiments \cite{lopezdeprado2016hrp}.

Two structural observations drive this thesis. First, the dendrogram
exists only to produce an ordering: allocation itself (step iii) consumes an
ordered list plus cluster variances, not the tree. Any procedure that yields
an economically meaningful asset ordering can therefore replace the
clustering step wholesale; this observation motivates causal-ordered
bisection (Section~\ref{sec:d2}). Second, every input to steps
(i)--(ii) is symmetric by necessity, not by accident: linkage clustering
merges clusters by comparing distances, and a distance is symmetric by
definition, so the pipeline is structurally blind to directionality. If the
true data-generating process contains one-way influence ($A$ moves $B$ with
a sign and magnitude that $B$ does not reciprocate), no correlation-fed
hierarchical allocator can represent that fact. This is the interface deficit of
Section~\ref{sec:motivation}: not a flaw of correlation estimation but of
what the clustering step can accept. Making that interface accept directional
information, without breaking what makes HRP work, is challenge C2.

The limitation is family-wide, not HRP-specific. HRP seeded a family of
hierarchical allocators (hierarchical clustering-based asset allocation
and its equal-risk-contribution variant \cite{raffinot2017hierarchical},
and nested clustered optimisation \cite{lopezdeprado2020mlam}) whose members
vary the linkage rule, the risk measure or the within-cluster optimiser,
but every one of which consumes a symmetric similarity structure at the
same first step. A decomposition of what directed structure is worth at
that interface therefore speaks to the family, not to one algorithm; HRP
is used here as its simplest and most-studied representative, so that any
effect found is attributable to the input rather than to allocator
elaborations. The interface claim is structural; whether the
measured decomposition also transfers across the family's
tree-reading rules is an empirical question, tested on a second member
(a HERC variant \cite{raffinot2018herc}) in
Section~\ref{sec:decomposition}, with a mixed outcome reported there.

A separate and well-established route to better allocation changes not the
allocator but the quality of its second-moment input. Linear shrinkage pulls
the noisy sample covariance towards a structured target with an analytically
optimal intensity, trading a small bias for a large reduction in estimation
variance \cite{ledoitwolf2004}; removing a common market factor through a
single-factor model and retaining the residual (de-factored) covariance
suppresses the dominant source of comovement noise in equity panels. Both
constructions condition the covariance estimate without asserting any graph
over the assets. Chapter~\ref{chap:evaluation} uses exactly these two
estimators, Ledoit--Wolf shrinkage and a de-factored residual covariance, as
graph-free comparison controls for the mechanism behind the SEM-implied
covariance.

The same author's later programme supplies the evaluation standard. The core
claim of the backtest-overfitting agenda \cite{bailey2014deflated,
lopezdeprado2023causal} is that the hard problem in quantitative finance is
not finding a configuration that looks good in a backtest, but establishing
that its performance is not an artefact of the search that found it. Two
instruments follow. The Probabilistic Sharpe Ratio \cite{bailey2012psr}
replaces the Gaussian sampling theory of the Sharpe ratio, untenable at the
excess kurtosis of ${\approx}\,\rsKurtosis$ measured here, with a version
that corrects for the sample's skewness, kurtosis and length; the
Deflated Sharpe Ratio \cite{bailey2014deflated} then raises the
benchmark from zero to the best Sharpe expected from $N$ independent trials
under the null of no skill, directly discounting selection bias. For
family-wise comparison against a benchmark, White's Reality Check
\cite{white2000reality} and Hansen's SPA test \cite{hansen2005spa} bootstrap
the null that no candidate beats it, and the Model Confidence Set
\cite{hansen2011mcs} returns the subset of strategies statistically
indistinguishable from the best. This battery is applied to every claim in
Chapter~\ref{chap:evaluation}: the Deflated Sharpe Ratio deflates each
headline Sharpe against all \rsNtrials\ configurations this project
evaluated, while the SPA and MCS run over the smaller family
of candidates relevant to each question. Several headline-looking effects do
not survive it, and reporting that is treated as part of the result rather
than a failure.

\section{Theme II: causal graphs on equities, without allocation
(Howard, Lohre and Mudde)}
\label{sec:howard-theme}

Howard, Lohre and Mudde \cite{howard2025causal} is the closest prior work: a
large-scale application of DYNOTEARS to S\&P~500 constituents over
1989--2022 (${\approx}\,357$ names per estimated network), re-estimated on
a rolling basis, structurally the same discovery machinery as
Chapter~\ref{chap:discovery}. Three of their empirical findings shape this
study's design. First, lagged edges are negligible at daily-to-monthly
horizons: the contemporaneous block $W$ contains essentially all the
structure, which justifies this project's use of the contemporaneous
asset--asset block as ``the graph'' and the lag-1 truncation of
Section~\ref{sec:asset-graph}.
Second, causal networks carry regime information that correlation
misses: rising network density precedes lower returns and flags recessions,
and their centrality-based signal earns its alpha in stable rather than
stressed periods. This motivates the regime-conditional evaluation of
Section~\ref{sec:regimes}, where the orientation effect is found to echo
their calm-regime pattern. Third, their overall verdict is measured:
causal networks complement rather than consistently beat correlation-based
methods. The research question sharpens that claim into a decomposition, asking whether the
complement operates through the skeleton or through the orientations.

What they do not do defines the gap this project fills. Their three
use-cases (peer groups for factor neutralisation, a low-centrality factor,
and a density-based market-timing signal) all consume the graph as a
similarity structure or a scalar diagnostic. The graph is never handed to an
allocator: no portfolio weights are ever computed from the discovered
network, and the directionality of the edges is used only implicitly (via
centrality). Consequently their pipeline cannot ask, let alone answer, how
much of the graph's allocation value lies in its orientations. Tang
\cite{tang2024trading} builds trading signals from time-series causal
discovery (including VARLiNGAM) and finds computational cost the binding
constraint, a warning this project addressed with a content-keyed discovery
cache (Section~\ref{sec:cache}) rather than by shrinking the experiment.

\section{The causal-discovery toolkit}
\label{sec:cd-tools}

The pipeline needs three facts from the causal-discovery literature
\cite{spirtes2000causation,kaddour2022causalml}. (i) Score-based
discovery scales and accepts priors. Constraint-based methods (PC, FCI)
scale poorly with dimension and lean on faithfulness, an assumption that is
fragile in markets,
where offsetting mechanisms can leave causally-connected pairs with near-zero
marginal correlation \cite{howard2025causal}. NOTEARS \cite{zheng2018notears}
recasts acyclicity as the smooth constraint
$h(W)=\operatorname{tr}(e^{W\circ W})-d=0$, making DAG learning a continuous
optimisation; DYNOTEARS \cite{pamfil2020dynotears} extends it to the
structural VAR $X = XW + \sum_p Y_p A_p + Z$ with acyclicity on the
contemporaneous $W$ only, and natively supports forbidden-edge lists, which
is how the directional prior of C1 is enforced. It is the primary method
here.
(ii) Non-Gaussianity identifies direction. The LiNGAM family
\cite{shimizu2006lingam} exploits non-Gaussian residuals to identify a unique
causal ordering; VARLiNGAM \cite{hyvarinen2010varlingam} is its SVAR form and
serves as the robustness comparator: its identification rests on an
assumption (heavy-tailed residuals) that financial returns satisfy at short
windows but progressively lose at longer ones as aggregation pulls residuals
toward Gaussianity, a mechanism the results repeatedly exhibit.
(iii) A directed graph without trustworthy magnitudes exists in
closed form. A ridge-regularised VAR(1) coefficient matrix is a
Granger-causal \cite{granger1969investigating} directed graph: it carries a
defensible directional ordering, but its shrunken one-lag
coefficients are not structural effect sizes. That combination makes it a
natural control for dissociating what ``orientation'' contains: the
arrows themselves versus magnitudes fit to propagate shocks.

\section{Research gap}
\label{sec:gap}

Table~\ref{tab:gap} positions the work; the \HRP{} control of
Section~\ref{sec:question} is exactly the HRP row. Reading across the
columns: HRP fixes the allocator but is blind to direction at the interface;
Howard et al.\
discover directed graphs at scale but never allocate from them. Neither
sets hierarchical portfolio weights from a discovered asset--asset
causal graph, and consequently neither can decompose the graph's allocation
value into skeleton and orientation, which is the decomposition the research
question demands and
Chapters~\ref{chap:sensitivity}--\ref{chap:evaluation} deliver.

The novelty claim needs care, because allocating from a graph over assets is
not itself new. The correlation-network tradition builds portfolios from
minimum-spanning-tree and PMFG filtrations of the correlation matrix, most
directly in the finding that weight placed in the peripheries of the network
outperforms \cite{pozzi2013spread}, and in network-centrality portfolio
selection \cite{peralta2016network}; clustering assets on a graph and then
allocating down its structure descends from that tradition, and the skeleton
arm of this thesis is its conceptual descendant. Nor is directed-graph
allocation unprecedented: Granger-causality networks over financial
institutions are an established systemic-risk measurement tool
\cite{billio2012econometric}, and V\'yrost, Ly\'ocsa and Baum\"ohl set
portfolio weights directly from such directed networks \cite{vyrost2019network}.
What survives, to my knowledge, is a sharper residual gap: no prior work sets
hierarchical portfolio weights from a discovered structural, score-based
causal graph, as opposed to a network of pairwise predictive tests, and none
decomposes the allocation value of such a graph into its skeleton and
orientation halves. That decomposition is exactly what this thesis delivers.

\begin{table}[t]
  \centering
  \small
  \caption{Positioning against the closest related approaches. The HRP row is
  the graph-blind \HRP{} control of this study.}
  \label{tab:gap}
  \begin{tabular}{l l c c c c}
    \toprule
    Approach & Clusters on & Causal & Weights & Asymmetric & Skeleton vs \\
             &             & discovery? & from graph? & allocation? & orientation? \\
    \midrule
    HRP \cite{lopezdeprado2016hrp} ($\equiv$ \HRP{}) & correlation & no & n/a & no & n/a \\
    MST/PMFG networks \cite{pozzi2013spread,peralta2016network} & corr.\ network & no & yes & no & n/a \\
    Granger networks \cite{vyrost2019network} & (no clustering) & no & yes & yes & no \\
    Howard et al.\ \cite{howard2025causal} & (no allocation)  & yes & no  & no & no \\
    \textbf{This work}                     & \textbf{causal graph} & \textbf{yes} & \textbf{yes} & \textbf{yes} & \textbf{yes} \\
    \bottomrule
  \end{tabular}
\end{table}
